\documentclass[10pt,a4paper]{amsart}
\usepackage[margin=19mm]{geometry}
\usepackage{amsmath,amssymb,mathtools}
\usepackage[scr=rsfs,cal=euler]{mathalfa}
\usepackage{xfrac}
\usepackage[unicode,hidelinks,bookmarks=false]{hyperref}
\newcommand{\ie}{i.e.\ }
\newcommand{\cf}{cf.\ }
\newcommand{\resp}{respectively\ }
\newcommand{\Subsection}{\subsection}
\providecommand{\nobreakdash}{\nobreak\hbox{-}}
\setlength{\emergencystretch}{2em}
\multlinegap0pt
\usepackage{bm}
\usepackage{bbm}
\usepackage{dsfont}
\usepackage{xspace}
\usepackage{cancel}
\usepackage{mathtools}
\usepackage{amsthm}
\makeatletter
\@ifundefined{theorem}{\newtheorem{theorem}{Theorem}}{}
\@ifundefined{proposition}{\newtheorem{proposition}[theorem]{Proposition}}{}
\makeatother
\title[Critical Ledgers and Scale-Defect Cascades for Navier-Stokes]{Critical Ledgers and Scale-Defect Cascades for Navier-Stokes}
\author{Runlong Yu}
\date{Source: arXiv:2606.13887v1 (CC BY 4.0)}
\begin{document}
\maketitle

\noindent\emph{Source and licence.} Critical Ledgers and Scale-Defect Cascades for Navier-Stokes. Version arXiv:2606.13887v1 (CC BY 4.0), distributed under the Creative Commons Attribution 4.0 International licence, as recorded in the arXiv licence metadata for this exact version and shown on the abstract page https://arxiv.org/abs/2606.13887v1. Full rights notice: licenses/arxiv-2606.13887-CC-BY-4.0.txt.\par\smallskip

\noindent\textit{Editorial scope.} Counted unit: the Proposition whose source label is \texttt{prop:no-phantom-gap-equivalence} in the article (source0.tex lines 3279--3294, source label \texttt{prop:no-phantom-gap-equivalence}), delivered together with the article's own complete original proof of it (source0.tex lines 3296--3404, one \texttt{proof} environment of 109 lines). No mathematics is written, repaired or re-derived: statement and proof are the article's own text line for line. Required context retained uncounted: none. Every definition, display and label of those lines is retained unchanged. Omitted from this package and not redistributed: 1--3278, 3295--3295, 3405--5742. Nothing inside a retained excerpt is omitted or altered: every retained body is delivered exactly as the article's lines are written, so no omission receipt of any kind accompanies this package. Citations: the retained excerpts carry no citation command at all, so no bibliography entry of the article is transcribed and no reference is redistributed. Reference adaptation: none was needed and none was made; every reference inside the delivered unit resolves inside it, so no unresolved reference or citation is produced. Printed slips kept literal: no printed word, symbol, hypothesis or number of the counted statement or of its proof was altered. Layout and class: the article's own class is \texttt{article} with packages including geometry, inputenc, fontenc, lmodern, amsmath, amssymb, amsthm, mathtools, mathrsfs, enumitem, booktabs, hyperref, xcolor, microtype; the wrapper is the lane's pinned \texttt{amsart} editorial wrapper at 10pt on A4, which restates the theorem-like environments the retained text needs and adds the article's own macro definitions that the retained text needs (none), with extra wrapper packages (bm, bbm, dsfont, xspace, cancel, mathtools). The delivered numbering is the unit's own. Assets: no retained span contains a figure environment, a graphics-inclusion command, a PDF-page inclusion, a table or a TikZ picture; no image file is copied into this package, the package declares no asset dependency and no image file is redistributed with it. Licence: version arXiv:2606.13887v1 of this article is distributed under the Creative Commons Attribution 4.0 International licence, as recorded in the arXiv licence metadata for this exact version and shown on the abstract page https://arxiv.org/abs/2606.13887v1. A full rights notice is delivered at licenses/arxiv-2606.13887-CC-BY-4.0.txt. Characters: every retained excerpt is pure ASCII, so the delivered engine, XeLaTeX with its default Latin Modern text font, reports no missing character.
\begin{proposition}[No exact dynamic phantom is equivalent to positive dynamic gap]
\label{prop:no-phantom-gap-equivalence}
Assume that \(\mathcal D_\Lambda/\operatorname{Im}G_\Lambda\) is
finite-dimensional, that the combined defect-size functional descends
continuously and positively one-homogeneously to the quotient, and that the
zero set defining \(\mathcal Z_\Lambda\) is conic.  Then
\[
        \mathcal Z_\Lambda
        \subset
        \operatorname{Im}G_\Lambda
\]
if and only if
\[
        \mu_\Lambda^{\rm dyn}>0.
\]
\end{proposition}

\begin{proof}
Suppose first that
\[
        \mu_\Lambda^{\rm dyn}=0.
\]
Then there exists a sequence
\(\mathfrak D_\Lambda^{(n)}\in\mathcal D_\Lambda\) such that
\[
        \operatorname{dist}_\Lambda
        (\mathfrak D_\Lambda^{(n)},\operatorname{Im}G_\Lambda)=1
\]
and
\[
        \|O_\Lambda\mathfrak D_\Lambda^{(n)}\|
        +
        \|\mathcal E_\Lambda(\mathfrak D_\Lambda^{(n)})\|
        +
        \operatorname{Rep}_\Lambda(\mathfrak D_\Lambda^{(n)})
        +
        [\mathsf P_\Lambda(\mathfrak D_\Lambda^{(n)})]_+
        \longrightarrow 0 .
\]
Passing to the quotient unit sphere and using finite-dimensional
compactness, we may extract a subsequence converging to a class
\[
        [\mathfrak D_\Lambda^\ast]
        \in
        \mathcal D_\Lambda/\operatorname{Im}G_\Lambda
\]
with
\[
        \|[\mathfrak D_\Lambda^\ast]\|=1.
\]
By continuity,
\[
        O_\Lambda\mathfrak D_\Lambda^\ast=0,
        \qquad
        \mathcal E_\Lambda(\mathfrak D_\Lambda^\ast)=0,
        \qquad
        \operatorname{Rep}_\Lambda(\mathfrak D_\Lambda^\ast)=0,
        \qquad
        [\mathsf P_\Lambda(\mathfrak D_\Lambda^\ast)]_+=0.
\]
Thus
\[
        \mathfrak D_\Lambda^\ast\in\mathcal Z_\Lambda.
\]
If
\[
        \mathcal Z_\Lambda\subset\operatorname{Im}G_\Lambda,
\]
then
\[
        [\mathfrak D_\Lambda^\ast]=0
\]
in the quotient, contradicting
\[
        \|[\mathfrak D_\Lambda^\ast]\|=1.
\]
Hence \(\mu_\Lambda^{\rm dyn}>0\).

Conversely, suppose that
\[
        \mathcal Z_\Lambda
        \not\subset
        \operatorname{Im}G_\Lambda.
\]
Then there exists
\[
        \mathfrak D_\Lambda^\ast\in\mathcal Z_\Lambda
\]
with
\[
        \operatorname{dist}_\Lambda
        (\mathfrak D_\Lambda^\ast,\operatorname{Im}G_\Lambda)>0.
\]
Normalize it by setting
\[
        \widetilde{\mathfrak D}_\Lambda
        :=
        \frac{\mathfrak D_\Lambda^\ast}
        {\operatorname{dist}_\Lambda
        (\mathfrak D_\Lambda^\ast,\operatorname{Im}G_\Lambda)} .
\]
Since \(\mathcal Z_\Lambda\) is conic,
\[
        \widetilde{\mathfrak D}_\Lambda\in\mathcal Z_\Lambda,
\]
and therefore
\[
        \operatorname{dist}_\Lambda
        (\widetilde{\mathfrak D}_\Lambda,\operatorname{Im}G_\Lambda)=1,
\]
while
\[
        O_\Lambda\widetilde{\mathfrak D}_\Lambda=0,
        \qquad
        \mathcal E_\Lambda(\widetilde{\mathfrak D}_\Lambda)=0,
        \qquad
        \operatorname{Rep}_\Lambda(\widetilde{\mathfrak D}_\Lambda)=0,
        \qquad
        [\mathsf P_\Lambda(\widetilde{\mathfrak D}_\Lambda)]_+=0.
\]
Therefore
\[
        \mu_\Lambda^{\rm dyn}=0.
\]
This proves the equivalence.
\end{proof}

\end{document}
